\chapter{Results}
\label{ch:results}

This chapter reports the experiments carried out on the prototype and uses them to answer the research questions. ChainBoard was evaluated in four ways. First, every function was exercised on a real Fabric network and the permission rules were tested with every role. Then the audit verification was tested on untouched histories and on reports manipulated on purpose. The speed of the system was measured with a benchmark. Finally, ten people used the prototype in a small user study. \Cref{tab:eval-overview} shows which research question each experiment answers.

\begin{table}[htbp]
    \centering
    \small
    \begin{tabular}{p{1.1cm}p{5.6cm}p{5.4cm}}
        \toprule
        RQ & What was evaluated & Experiment \\
        \midrule
        RQ1 & The project management functions & Scenarios on the real network, automated tests \\
        RQ2 & Permissions enforced by the chaincode & Every action tried with every role \\
        RQ3 & Verifiable audit trail & Untouched and manipulated histories \\
        RQ4 & Cost and limits & Benchmark, threat model \\
        RQ1, RQ3 & Whether people can use the board and understand the history check & User study with ten participants \\
        \bottomrule
    \end{tabular}
    \caption{Overview of the evaluation.}
    \label{tab:eval-overview}
\end{table}

\section{Setup}
\label{sec:results-setup}

All experiments were run on one machine hosting the Fabric test network (two organizations with one peer each, one orderer and one channel), a local IPFS node, the backend and the frontend. \Cref{tab:environment} lists the details. The machine is an ordinary laptop, not a server, which matters when reading the performance figures. The demonstration data was loaded with a seed script that creates five accounts (an owner, an administrator, two contributors and an outsider), four projects and their tasks, all through the normal API, so that every record is a real transaction.

\begin{table}[htbp]
    \centering
    \small
    \begin{tabular}{ll}
        \toprule
        Item & Value \\
        \midrule
        CPU & Intel Core i5-1335U (13th gen), 12 logical CPUs \\
        RAM available to WSL & 7.6\,GiB \\
        Operating system & Windows 11, WSL2 with Ubuntu 24.04.1 LTS \\
        Docker, Node.js & 26.0.0, 20.20.2 \\
        Hyperledger Fabric & 2.5.16 \\
        Peers, orderers & 2 peers (one per organization), 1 orderer \\
        Batch timeout of the orderer & 2\,s (test network default) \\
        IPFS (Kubo) & 0.29.0 \\
        \bottomrule
    \end{tabular}
    \caption{Test environment.}
    \label{tab:environment}
\end{table}

\section{Functional Evaluation}
\label{sec:results-functional}

The chaincode functions were first called directly with the Fabric peer command-line tool, without the backend, and the same steps were then repeated through the web interface. \Cref{tab:scenarios} lists the scenarios and their outcome. All of them behaved as expected, and the state reached through the interface was the same as the state reached with the command-line tool, so the backend introduces no differences of its own.

\begin{table}[htbp]
    \centering
    \small
    \begin{tabular}{p{0.6cm}p{5.4cm}p{5.0cm}p{0.9cm}}
        \toprule
        \# & Scenario & Observed & Result \\
        \midrule
        F1 & Create a project & Owner is the only member, status \texttt{active} & pass \\
        F2 & Add a member with a role & Member appears in the project with that role & pass \\
        F3 & Create a task & Status \texttt{todo}, no assignee & pass \\
        F4 & Assign the task to a member & Assignee stored; a non-member is refused & pass \\
        F5 & Move the task \texttt{todo}\,$\rightarrow$\,\texttt{in-progress}\,$\rightarrow$\,\texttt{done} & Each move stored and visible on the board & pass \\
        F6 & Jump \texttt{todo}\,$\rightarrow$\,\texttt{done} & Refused as an invalid transition & pass \\
        F7 & Comment on the task & Comment stored with its author & pass \\
        F8 & Attach a file & File uploaded to IPFS, CID stored on the ledger, same file returned by the gateway & pass \\
        F9 & Archive the task and the project & Removed from the board; archived project accepts no new tasks & pass \\
        F10 & Read the history of the task & Every version, oldest first, with author and time & pass \\
        F11 & Sign up, sign in, change password & Account events stored as transactions & pass \\
        \bottomrule
    \end{tabular}
    \caption{Functional scenarios, run on the real network.}
    \label{tab:scenarios}
\end{table}

\section{Permissions}
\label{sec:results-permissions}

To check the permission rules, each action was tried with each of the five seeded accounts, both through the interface and directly against the API, where the interface cannot hide a button. \Cref{tab:perm-results} shows the outcome. Every forbidden action was refused by the chaincode with ``Permission denied'', which the backend returns as HTTP 403, and nothing was written to the ledger. Every allowed action succeeded.

\begin{table}[htbp]
    \centering
    \small
    \begin{tabular}{p{4.6cm}cccc}
        \toprule
        Action & Owner & Admin & Contributor & Outsider \\
        \midrule
        Add a contributor & yes & yes & 403 & 403 \\
        Add an admin & yes & 403 & 403 & 403 \\
        Archive the project & yes & 403 & 403 & 403 \\
        Create a task, comment & yes & yes & yes & 403 \\
        Take a free task & yes & yes & yes & 403 \\
        Assign a task to someone else & yes & yes & 403 & 403 \\
        Change a task assigned to another person & yes & yes & 403 & 403 \\
        Archive a task & yes & yes & 403 & 403 \\
        \bottomrule
    \end{tabular}
    \caption{Outcome of each action for each role (yes = allowed and committed, 403 = refused by the chaincode).}
    \label{tab:perm-results}
\end{table}

Two more checks concern the acting person. A request whose body names a different person is still recorded under the signed-in user, because the backend ignores that field. A call that reaches the chaincode without any acting person is refused by the chaincode itself. Refused calls leave nothing in the history. This is by design: the history records changes, not attempts.

\section{Audit Trail Verification}
\label{sec:results-audit}

The audit verification is the main contribution of the thesis, so it was tested in two situations: on real histories that nobody had touched, and on downloaded reports that were changed on purpose.

On untouched data, ``Verify audit trail'' was run on every task and project of the demonstration data. Each time, all records were found in committed blocks, and the offline verifier computed from the downloaded file the same head digest that the application had shown.

\Cref{tab:tamper} shows the manipulations. In every case the change was detected, and the verifier named the first record that no longer matched. The last row matters too: when the peer could not be reached, the records were marked ``not checked'' instead of ``missing'', so a network problem is not mistaken for tampering.

\begin{table}[htbp]
    \centering
    \small
    \begin{tabular}{p{4.6cm}p{2.0cm}p{5.0cm}}
        \toprule
        Change made to the report or history & Detected & What the user sees \\
        \midrule
        None (untouched history) & -- & Verified, same head digest online and offline \\
        Edit one field of one record & yes & Digest mismatch from that record on \\
        Delete a record & yes & Digest mismatch from the gap on \\
        Swap two records & yes & Digest mismatch from the first swapped record \\
        Insert a fake record & yes & Digest mismatch at the fake record, every later record out of order, record count wrong \\
        Transaction unknown to the peer & yes & Record reported as missing \\
        Peer not reachable & -- & Records reported as not checked \\
        \bottomrule
    \end{tabular}
    \caption{Manipulations of an audit report and the result of the verification.}
    \label{tab:tamper}
\end{table}

As a concrete example, the report of a demonstration task with four records was downloaded, and the offline verifier accepted it unchanged. A copy of the first record with a different transaction identifier was then inserted into the file. The verifier rejected the modified file: the digest of the fake record did not match, every record after it was out of order, and the record count no longer matched.

What the check cannot show is also important. It proves that the history was not changed after it was written. It does not prove that the person recorded in it was really the person who acted, because the backend reports who is acting.

\section{Automated Tests}
\label{sec:results-tests}

The backend tests run the real chaincode on an in-memory implementation of the Fabric interface, behind the real Express application, so they need no network. The frontend tests check the decision logic of each screen without a browser. \Cref{tab:tests} gives the numbers; all 278 tests pass.

\begin{table}[htbp]
    \centering
    \small
    \begin{tabular}{p{7.4cm}rr}
        \toprule
        Area & Tests & Passed \\
        \midrule
        Projects, tasks, transitions, ownership from the session & 18 & 18 \\
        Permissions (every role and action) & 14 & 14 \\
        History, digest chain and audit verification & 17 & 17 \\
        Accounts: sign-up, sign-in, sessions, passwords, Google, e-mail & 98 & 98 \\
        Security: CSRF, headers, route protection, rate limits, bad input & 24 & 24 \\
        Benchmark statistics helpers & 5 & 5 \\
        \midrule
        Backend total & 176 & 176 \\
        Frontend decision logic (14 files) & 102 & 102 \\
        \midrule
        Total & 278 & 278 \\
        \bottomrule
    \end{tabular}
    \caption{Automated tests (Node.js test runner; backend about 8\,s, frontend under 1\,s).}
    \label{tab:tests}
\end{table}

Tests on a simulated ledger are useful for the business rules, but they do not replace running on the network: the order and the field names of the history returned by the real Fabric library differ from those of the in-memory implementation, and this difference became visible only on the real network.

\section{Performance}
\label{sec:results-performance}

The benchmark script sends real requests through the backend, so every write goes through endorsement, ordering and commit. It uses its own temporary account and project. It measures single writes (40 \texttt{createTask} calls one after the other), write throughput with 1, 2, 4 and 8 concurrent clients, contention (8 clients editing the same task at once) and reads (200 reads of one task and of all tasks of a project). The benchmark was run three times in a row, and the tables report, for each value, the median of the three runs. The runs differed very little: write latencies by less than 2\%, reads by a few milliseconds.

\begin{table}[htbp]
    \centering
    \small
    \begin{tabular}{lrrr}
        \toprule
        Operation & p50 (ms) & p95 (ms) & Failed \\
        \midrule
        Write one task (\texttt{createTask}) & 2087 & 2129 & 0 / 40 \\
        Read one task & 30 & 44 & 0 / 200 \\
        Read all tasks of a project & 30 & 57 & 0 / 200 \\
        \bottomrule
    \end{tabular}
    \caption{Latency of single operations through the whole stack (median of three runs).}
    \label{tab:benchmark}
\end{table}

\begin{table}[htbp]
    \centering
    \small
    \begin{tabular}{rrrr}
        \toprule
        Concurrent clients & Throughput (tx/s) & p50 (ms) & p95 (ms) \\
        \midrule
        1 & 0.48 & 2076 & 2099 \\
        2 & 0.95 & 2101 & 2124 \\
        4 & 1.89 & 2116 & 2142 \\
        8 & 3.73 & 2143 & 2172 \\
        \bottomrule
    \end{tabular}
    \caption{Write throughput with concurrent clients writing different tasks.}
    \label{tab:throughput}
\end{table}

\paragraph*{Writes and reads.} A write takes about 2.1 seconds, and the spread is very small: the slowest of 40 writes took 2.16 seconds. Almost all of this time is waiting. When there are only a few transactions, the orderer of the test network cuts a block every 2 seconds, so a write costs the batch timeout plus about 100\,ms for endorsement, validation and commit. Reads take about 30\,ms, around 70 times faster, because one peer answers from its local state without waiting for a block. With the small amount of data used here, reading all tasks of a project costs the same as reading one task.

\paragraph*{Throughput.} Throughput doubles every time the number of clients doubles, from 0.48 to 3.73 transactions per second, while the latency of each write stays almost the same (from 2.08 to 2.14 seconds). The reason is that all writes sent within the same 2-second window end up in the same block. With 8 clients the system had not reached its limit yet; at this load, the block timeout, not the processor, decides the speed.

\paragraph*{Contention.} When 8 clients edited the same task at once, only 1 edit was committed and the other 7 were rejected with an MVCC read conflict, in all three runs. The backend returns these as HTTP 409 so that the user can try again. This is the correct behavior of Fabric: only one change to a key can win in each block. It is also the reason why the data model gives every task, project and account its own key, so that people working on different tasks never block each other.

\paragraph*{How to read these numbers.} Everything ran on one machine, with the load generator sharing the processor with the network, so the figures show orders of magnitude, not the limits of Fabric. They are in line with the factors identified by \textcite{thakkar2018performance}, mainly the block settings of the orderer and conflicts on the same key. A conventional database behind the same API would be faster; the extra time is the price of a history that no single administrator can rewrite. The 2-second wait could be shortened by lowering the batch timeout, at the cost of more and smaller blocks. For a team tool, where a person changes a task every few minutes, a write time of about 2 seconds is noticeable but acceptable, especially because the interface updates the card immediately and confirms in the background.

% NOT included in main.tex yet. After the user test, paste this section back
% into chapters/results.tex, just before \section{Threat Model}.

\section{User Study}
\label{sec:results-usability}

To find out whether people can actually use ChainBoard and understand its audit check, I ran a small usability study.

\paragraph{Participants and procedure.} \textit{tbd} people took part (\textit{tbd} students and \textit{tbd} developers; \textit{tbd} had used Trello or Jira before, \textit{tbd} had used a blockchain application). Each session took about 25 minutes on my laptop. Participants followed a printed sheet of 12 tasks in a prepared project (\Cref{tab:usability-tasks}) and were asked to think aloud. I did not help unless someone was stuck for more than two minutes, which I counted as a failure. Afterwards they filled in the System Usability Scale (SUS)~\cite{brooke1996sus} and two extra statements about the audit trail, on the same 1--5 scale, plus three open questions.

\begin{table}[tb]
    \centering
    \small
    \begin{tabular}{p{0.5cm}p{6.6cm}rr}
        \toprule
        \# & Task & Success & Median time (s) \\
        \midrule
        T1 & Sign in & \textit{tbd}/\textit{n} & \textit{tbd} \\
        T2 & Open the project and its board & \textit{tbd}/\textit{n} & \textit{tbd} \\
        T3 & Find the overdue task & \textit{tbd}/\textit{n} & \textit{tbd} \\
        T4 & Assign it to yourself & \textit{tbd}/\textit{n} & \textit{tbd} \\
        T5 & Move it to In progress & \textit{tbd}/\textit{n} & \textit{tbd} \\
        T6 & Add a comment & \textit{tbd}/\textit{n} & \textit{tbd} \\
        T7 & Create a task with priority and due date & \textit{tbd}/\textit{n} & \textit{tbd} \\
        T8 & Attach a file & \textit{tbd}/\textit{n} & \textit{tbd} \\
        T9 & Mark the task Done & \textit{tbd}/\textit{n} & \textit{tbd} \\
        T10 & Find who completed another task, and when & \textit{tbd}/\textit{n} & \textit{tbd} \\
        T11 & Verify the audit trail and explain the result & \textit{tbd}/\textit{n} & \textit{tbd} \\
        T12 & Try to change another person's task, explain why it fails & \textit{tbd}/\textit{n} & \textit{tbd} \\
        \bottomrule
    \end{tabular}
    \caption{Tasks of the user study, success rate and median completion time.}
    \label{tab:usability-tasks}
\end{table}

\paragraph{Task success.} \ag{Two or three sentences: overall success rate, which tasks were hardest and why (use what people said while thinking aloud).} For T11 and T12 a task counted as successful only if the participant could also explain the result in their own words, for example that the history is stored on a ledger and cannot be changed, or that only the assignee may change a task.

\paragraph{SUS score.} Each SUS answer was converted to a 0--100 score in the standard way: for odd items subtract 1 from the answer, for even items subtract the answer from 5, add the ten values and multiply by 2.5. \Cref{tab:sus} gives the results. An average SUS score is around 68~\cite{lewis2018sus}, so \ag{compare the mean with 68 in one sentence}.

\begin{table}[tb]
    \centering
    \small
    \begin{tabular}{lrrrr}
        \toprule
        Measure & Mean & SD & Min & Max \\
        \midrule
        SUS score (0--100) & \textit{tbd} & \textit{tbd} & \textit{tbd} & \textit{tbd} \\
        ``I understood what the result of Verify audit trail means'' (1--5) & \textit{tbd} & \textit{tbd} & \textit{tbd} & \textit{tbd} \\
        ``Being able to check a task's history makes me trust it more'' (1--5) & \textit{tbd} & \textit{tbd} & \textit{tbd} & \textit{tbd} \\
        \bottomrule
    \end{tabular}
    \caption{Questionnaire results ($n = \textit{tbd}$).}
    \label{tab:sus}
\end{table}

\paragraph{What participants said.} \ag{Summarize the open answers and think-aloud notes in a short paragraph or a small table: the two or three most common problems, what people liked, and one or two short quotes.}

\paragraph*{Changes made after the first session.} The study was used to improve the prototype while it was running. In the first session, the participant did not find how to set the priority and the due date of a new task (T7): the create-task bar showed only a title field and an ``Add task'' button, and the other fields were hidden behind a small ``More options'' link. After this session the create-task card was redesigned. The board now shows a single ``Add new task'' button, which opens a card that has a heading, a labelled title field, the priority as three visible buttons (low, medium, high), a visible due-date field, an optional description, and a ``Create task'' button with a hint when the title is missing; the task identifier is always generated automatically and is no longer shown. The result of the audit verification was also rewritten in plain language. The button ``Verify audit trail'' became ``Check that this history is genuine'', with one sentence explaining what it does; the result says, for example, ``This history is genuine: all 3 changes were found in the blockchain exactly as they are shown here'', the head digest is shown as a ``history code'' with a short explanation, and transaction identifiers moved under ``Technical details''. The following sessions used the new version. \ag{Say how many sessions used the old and the new version, and whether T7 became easier.}

\paragraph*{Changes made after the second session.} \ag{Start with what you observed in session 2 that led to these changes, e.g. where the participant got lost or hesitated.} After the second session the layout of the application was reorganized. A sidebar now lists the person's projects, so switching project takes one click from anywhere, and each project has a header with two tabs, ``Board'' and ``Members and history'', instead of separate pages and a ``Project settings'' button. A task no longer opens on its own page: it opens in a window over the board, with the status as three buttons under the title, a single line with priority, due date, assignee and task identifier, and below it the description, files, comments and activity. The separate ``Edit'' mode was removed; title, description, priority and due date are now changed in place and each change is saved as its own transaction. When a person tries to change something they are not allowed to, for example a task assigned to someone else, or to archive a task without being owner or admin, the interface now explains why instead of silently ignoring the click; the chaincode still enforces the same rules. Finally, the light theme was changed from grey to a warm off-white with near-black text and all text colours at or above the WCAG AA contrast of 4.5:1~\ag{cite WCAG 2.1}. \ag{Say which sessions used this version.}

With \textit{tbd} participants the study cannot give precise numbers, and most participants were students, not professional teams. It does show whether the main ideas, a familiar board and a visible audit check, make sense to people who did not build the system.


\section{Threat Model}
\label{sec:results-threats}

\paragraph*{Trust assumptions.} No single party is assumed to control both endorsement and ordering in the Fabric network. The backend is trusted, because it holds the Fabric identity of the organization, verifies the login session and tells the chaincode who is acting. Browsers are not trusted.

\paragraph*{What the design protects against.} Under these assumptions, an administrator cannot quietly change a status, a due date or a comment, because every change is a new transaction and the old versions stay in the history. A history changed in a copy of the data is detected by the digest chain, also offline. Two accounts cannot share a username, phone number or Google account. A user cannot perform an action that their role forbids, even if the backend has a bug, because the chaincode checks the role on every write.

\paragraph*{What it does not protect against.} \Cref{tab:threats} lists the remaining risks, which set the limits of the claims made in this thesis.

\begin{table}[htbp]
    \centering
    \small
    \begin{tabular}{p{3.4cm}p{8.6cm}}
        \toprule
        Threat & Why it remains, and how it could be fixed \\
        \midrule
        Compromised backend & Fabric only sees the organization, so the chaincode trusts the backend about who is acting. Fix: one Fabric certificate per person. \\
        One organization controls all endorsers & The test network has two organizations but one trust domain. Fix: require both organizations to endorse critical actions. \\
        Stolen password or session & Usual web risk, reduced by scrypt hashing, HttpOnly cookies, rate limits and CSRF checks. \\
        Reading without permission & Reads are not checked by role. Fix: role-checked reads or Fabric private data. \\
        Lost attachments & The CID proves what a file contained, but nobody is obliged to keep the file. Fix: pinning. \\
        Personal data on the ledger & Ledger data cannot be erased, which conflicts with the GDPR right to erasure~\cite{gdpr}. Fix: keep personal data off-chain and store only hashes. \\
        \bottomrule
    \end{tabular}
    \caption{Remaining threats.}
    \label{tab:threats}
\end{table}

\section{Support for Scrum and Kanban}
\label{sec:results-agile}

ChainBoard is in practice a Kanban board with a verifiable history. It does not implement Scrum; this was left out on purpose, because the research questions concern the integrity of the record, not planning features. \Cref{tab:agile} shows what is covered.

\begin{table}[htbp]
    \centering
    \small
    \begin{tabular}{p{4.4cm}p{1.6cm}p{6.0cm}}
        \toprule
        Element & Supported & Comment \\
        \midrule
        Board with workflow columns & yes & To do, in progress, done \\
        Tasks with priority, due date, assignee, comments, files & yes & \\
        Product backlog & partly & A flat list of tasks, no ordering \\
        Sprints, sprint goal & no & \\
        Scrum roles & no & Owner, admin and contributor only control permissions \\
        Estimation, burndown, velocity & no & \\
        Work-in-progress limits & no & \\
        Verifiable history of every change & yes & Not offered by Trello, Jira or similar tools \\
        \bottomrule
    \end{tabular}
    \caption{Scrum and Kanban elements in ChainBoard.}
    \label{tab:agile}
\end{table}

The history is still useful for Scrum events that the tool does not support directly. In a sprint review or retrospective, the team can see exactly when each task moved and who moved it, and a task that was reopened cannot be hidden.

\section{Answers to the Research Questions}
\label{sec:rq-answers}

\textbf{RQ1.} Yes. The Fabric-plus-IPFS pattern carries over to software project management with a different data model (users, projects, tasks, comments and attachments, with keys designed for a store that has no queries by field value) and a familiar board workflow. All eleven functional scenarios of \Cref{tab:scenarios} passed on the real network, through the command-line tool and through the web interface.

\textbf{RQ2.} Yes, with one qualification. In the permission experiment every forbidden action was refused by the chaincode, also when the API was called directly and the interface was bypassed, and refused calls left no trace on the ledger (\Cref{tab:perm-results}). The qualification is that the chaincode learns \emph{who} is acting from the backend, because Fabric authenticates the organization and not the person. Closing this gap requires one Fabric identity per person.

\textbf{RQ3.} The system checks every transaction against the peer, computes a digest chain over the history, and lets the user download a report that a stand-alone script verifies offline. Every manipulation tried in \Cref{tab:tamper} was detected, and an unreachable peer was never reported as tampering. The check shows that a history was not altered; it does not show that the person recorded as acting was really the one who acted.

\textbf{RQ4.} A write costs about 2.1 seconds, almost all of it waiting for the orderer to cut a block, while a read costs about 30\,ms. Throughput grew linearly from 0.48 to 3.73 transactions per second between 1 and 8 clients, and of 8 simultaneous edits of the same task only one was committed. The limits are those of \Cref{tab:threats}: identity, confidentiality of reads, permanence of personal data, availability of attachments, and the scale of the evaluation.

\section{Discussion}
\label{sec:results-discussion}

The experiments give a consistent picture. The functions that a software team expects from a board are available and behave the same through the interface and directly on the chaincode; the permission rules hold even when the interface is bypassed; and every manipulation of a history was detected, also offline. The price is speed: a write takes about two seconds instead of the few milliseconds of a conventional database, because it waits for a block. For the intended use, a team tool where each person changes a task every few minutes, this cost is acceptable.

Two observations go beyond the numbers. A test setup that simulates the ledger can differ from the real library in ways that hide defects, so tests on the real network remain necessary. And the most important limitation is architectural: because the Fabric gateway identifies an organization and not a person, a rule such as ``only an owner may archive a project'' holds only as far as the backend that reports who is acting can be trusted. The evaluation also used a local network with two organizations, not a geographically distributed or production-scale deployment; \Cref{ch:conclusion} discusses the consequences as limitations and future work.
